\small
\textbf{Abstract.}
Retailers plan store orders, promotions, category reviews and supplier terms
from the same demand at different grains. This paper asks whether one forecast, made for every item in
every store and summed to each coarser grain, can serve all of these decisions,
and which model does it best. It compares gradient-boosted trees,
pretrained foundation models, the Temporal Fusion Transformer and a world model that carries each item's state
forward under the price and promotion plan, on four public data sets (M5, six
Dominick's categories, Breakfast at the Frat and Corporaci\'on Favorita). The
world model and LightGBM are tuned by the same search, and every result is
scored on windows the search never saw. On items and stores seen in training,
the tuned world model and LightGBM are close, and their average is the
most accurate choice: it is in the 90\% model confidence set at \rqMcsAverage{}
of \rqMcsCells{} levels and the most accurate on the fast lines and the slow tail
on \rqAvgBestFast{} of the nine data sets. On items kept out of training, the
world model loses at most \rqCostWmItemHi{} in scaled error, against up to
\rqCostLgbItemHi{} for LightGBM. Settings tuned on one retailer carry to another,
and pooling other retailers' data improved neither model by more than
\rqPoolGainMax{}. On these data, the rule that follows is their average summed to every grain,
with the world model alone for items not yet trained on. Adopters should rerun
every comparison on their own data.

\vspace{4pt}
\textbf{Keywords:} retail forecasting; forecast grain; hierarchical evaluation;
forecast combination; promotions; world models; foundation models.
